\documentclass[10.5pt,a4paper]{article}
\usepackage[a4paper,margin=2.2cm,top=2.4cm,bottom=2.4cm]{geometry}
\usepackage{fontspec}
\usepackage{amsmath,amssymb,mathtools,bm}
\setmainfont{Noto Serif CJK SC}[BoldFont={Noto Serif CJK SC Bold}]
\setsansfont{Noto Sans CJK SC}[BoldFont={Noto Sans CJK SC Bold}]
\XeTeXlinebreaklocale "zh"
\XeTeXlinebreakskip = 0pt plus 1pt minus 0.1pt
\linespread{1.18}
\usepackage{xcolor}
\definecolor{accent}{HTML}{1F5FA8}
\definecolor{quiet}{HTML}{5A5A5A}
\definecolor{tint}{HTML}{EEF3FA}
\usepackage{titlesec}
\titleformat{\section}{\sffamily\bfseries\Large\color{accent}}{\thesection}{0.8em}{}
\titleformat{\subsection}{\sffamily\bfseries\large}{\thesubsection}{0.6em}{}
\titleformat{\subsubsection}{\sffamily\bfseries\normalsize}{\thesubsubsection}{0.5em}{}
\titlespacing*{\section}{0pt}{1.6em}{0.7em}
\titlespacing*{\subsection}{0pt}{1.2em}{0.5em}
\usepackage{enumitem}
\setlist{itemsep=0.25em,topsep=0.35em,leftmargin=1.6em}
\usepackage{array,booktabs,tabularx,xltabular,makecell}
\renewcommand{\arraystretch}{1.25}
\newcolumntype{L}{>{\raggedright\arraybackslash}X}
\newcolumntype{P}[1]{>{\raggedright\arraybackslash}p{#1}}
\usepackage{caption}
\captionsetup{font=small,labelfont=bf,labelsep=quad}
\renewcommand{\figurename}{图}
\renewcommand{\tablename}{表}
\renewcommand{\contentsname}{目录}
\usepackage{tikz}
\usetikzlibrary{positioning,arrows.meta,calc,fit}
\usepackage{fancyhdr}
\pagestyle{fancy}\fancyhf{}
\fancyhead[L]{\small\color{quiet}轨迹复形理论补全（未成型部分）}
\fancyhead[R]{\small\color{quiet}2026 年 10 月 15 日}
\fancyfoot[C]{\small\thepage}
\renewcommand{\headrulewidth}{0.4pt}
\usepackage[colorlinks=true,linkcolor=accent,citecolor=accent,urlcolor=accent]{hyperref}
\newcommand{\AFH}{AFRC$_{\mathrm H}$}
\newcommand{\AFE}{AFRC$_{\mathrm E}$}
\newcommand{\lnJ}{$\ln J$}
\newcommand{\keyp}[1]{\textbf{\sffamily #1}}
\newenvironment{finding}{\begin{center}\begin{minipage}{0.94\linewidth}\setlength{\parskip}{0.3em}\color{black}}{\end{minipage}\end{center}}
\setlength{\parskip}{0.35em}
\setlength{\parindent}{0pt}
\tolerance=2000 \emergencystretch=3em
\usepackage{amsthm}
\theoremstyle{definition}
\newtheorem{definition}{定义}[section]
\theoremstyle{plain}
\newtheorem{proposition}[definition]{命题}
\newtheorem{corollary}[definition]{推论}
\theoremstyle{remark}
\newtheorem{remark}[definition]{注}
\renewcommand{\proofname}{证明}
\newcommand{\status}[1]{\hfill{\small\color{accent}\textsf{[#1]}}}
\newcommand{\re}{\operatorname{Re}}
\newcommand{\im}{\operatorname{Im}}
\begin{document}
\begin{center}
{\sffamily\bfseries\LARGE 轨迹复形理论补全（未成型部分）}\\[0.5em]
{\sffamily\large 三点映射的复表示、合并穿越的锚定规则与物质胞腔复形}\\[0.8em]
{\color{quiet}2026 年 10 月 15 日}
\end{center}

\vspace{0.4em}
\noindent\fcolorbox{accent!40}{tint}{\begin{minipage}{0.96\linewidth}\vspace{0.4em}
\keyp{范围与结论}\quad
本稿不涉及已成型的论文（AFRC 两通道），只补统一理论中尚未成型的部分：物质三角形层、离散相对转角与 LAVD、激波签名与锚定规则、物质胞腔复形与空心胞腔、Hodge–Dirac 层的一般内积形式、阈值消除，以及与输运和掺混层的接口。

补全的核心是一个工具：二维三点仿射映射的复表示 $F\zeta=p\zeta+q\bar\zeta$。面积比、转角、各向异性（Beltrami 系数 $\mu=q/p$）和切向拉伸都可以由 $(p,q)$ 精确读出，合成律是封闭的，参考系变换下的行为也可以精确写出。由此得到四个新结果：
\begin{enumerate}[leftmargin=1.4em,itemsep=0.1em]
\item 离散相对转角对刚体参考系变换严格客观，并在光滑区收敛到 Haller 的相对动态转角；
\item 激波伪转动的来源是映射合成的非交换余项，上界为 $\sum_k\arcsin(|\mu_k||M_{k-1}|)$；把 $|\mu|$ 超阈的步合并为一个块之后，平面激波的伪转角从 0.02–0.2 rad 降到 $\le5\times10^{-3}$ rad；
\item 位置噪声在固定物质三角形上望远相消，在逐帧重剖时按 $\sqrt n$ 累积，所以“逐帧重剖”是现有实验中噪声的主要来源；
\item 单轴压缩满足 $|\mu|=\tanh(\tfrac12|\ln J|)$，给出只凭位置、只看单个三角形的激波签名，可以把 C 通道从“前提”降为“报告量”。
\end{enumerate}
每条命题标注状态：\textsf{[已证]}、\textsf{[已证+数值]}、\textsf{[数值]}、\textsf{[引用]} 或 \textsf{[待检验]}。
\vspace{0.4em}\end{minipage}}

\tableofcontents
\newpage

%=====================================================================
\section{约定与原则}

输入只有示踪点在快照上的位置 $x_i(t_s)$，$s=0,\dots,S-1$。物质元分三种：点对、物质三角形（以及第~\ref{sec:cells} 节的物质多边形）、物质区域。复平面记号为 $\zeta=X_1+iX_2$。分析区域记为 $U$，空间平均一律取面积加权，以与 Haller 的 $\bar\omega$ 定义一致。

已定稿的三条原则保持不变：拉伸取点对，转动与环量取物质三角形，收支取物质区域；跨激波只报告积分量的有限跳变；穿越步由时间上的合成或稳健统计吸收，不在空间上挖洞插补。本稿新增第四条：

\begin{finding}
\keyp{原则 4（锚定）。} 所有三角形层的量都在“块映射”上计算。块的切分由三角形自身的 Beltrami 系数决定：$|\mu|$ 超阈的连续步合并为一个块，块内不重锚；重锚只发生在块的边界，且只在形状门失效时进行。整个过程只用位置，只看单个三角形，不需要全局激波检测。
\end{finding}

%=====================================================================
\section{三点映射的复表示}\label{sec:complex}

\begin{definition}[复表示与 Beltrami 系数]
设物质三角形的三个顶点在锚定时刻 $t_a$ 位于 $X_a,X_b,X_c$，在时刻 $t$ 位于 $x_a,x_b,x_c$。三点唯一确定仿射映射 $x=FX+c$。在复平面上它写成
\begin{equation}
F\zeta=p\,\zeta+q\,\bar\zeta,\qquad
p=\tfrac12\big[(F_{11}+F_{22})+i(F_{21}-F_{12})\big],\qquad
q=\tfrac12\big[(F_{11}-F_{22})+i(F_{21}+F_{12})\big],
\end{equation}
称 $\mu=q/p$ 为该映射的 Beltrami 系数。
\end{definition}

\begin{proposition}[恒等式]\status{已证+数值}\label{prop:id}
对保向映射（$\det F>0$）：
\begin{enumerate}
\item 面积比 $J=\det F=|p|^2-|q|^2$；奇异值 $\sigma_{1,2}=|p|\pm|q|$；$|\mu|=(\sigma_1-\sigma_2)/(\sigma_1+\sigma_2)<1$。
\item 极分解 $F=RU$ 的转角为 $\arg p$。
\item 位移 $D=x-X$ 沿锚定三角形边界的环量满足 $c_t/A_t(t_a)=F_{21}-F_{12}=2\,\im p$，位移的散度满足 $\mathrm{tr}\,F-2=2(\re p-1)$。
\end{enumerate}
\end{proposition}
\begin{proof}
(1) $|p\zeta+q\bar\zeta|$ 在单位圆上的取值范围是 $[\,|p|-|q|,\,|p|+|q|\,]$，乘积即 $\det F$。(2) $F\zeta=e^{i\arg p}\big(|p|\zeta+q'\bar\zeta\big)$，其中 $q'=e^{-i\arg p}q$。后一映射的实矩阵为 $|p|I+\begin{psmallmatrix}\re q'&\im q'\\ \im q'&-\re q'\end{psmallmatrix}$，对称，且因 $|p|>|q'|$ 而正定，所以它就是 $U$，$R=e^{i\arg p}$。(3) 线性场在直边三角形上的线积分可由梯形公式精确求出，再用 Stokes 公式。
\end{proof}

\begin{proposition}[参考系变换]\status{已证+数值}\label{prop:frame}
设 $x=Q(t)y+b(t)$，$Q(t)$ 为转角 $\beta(t)$ 的旋转。块映射 $t_a\to t$ 在新参考系中的系数为
\begin{equation}
p\;\mapsto\;e^{-i(\beta(t)-\beta(t_a))}\,p,\qquad q\;\mapsto\;e^{-i(\beta(t)+\beta(t_a))}\,q .
\end{equation}
因此 $|p|$、$|q|$、$|\mu|$、$J$、$\sigma_{1,2}$ 都不变；$\arg p$ 对所有三角形平移同一个量 $-\Delta\beta$。
\end{proposition}
\begin{proof}
$y=Q^\top(x-b)$ 在复平面上即乘以 $e^{-i\beta}$。代入 $F$ 并整理 $\zeta$ 与 $\bar\zeta$ 的系数即得。数值上，随机映射与随机转角下 $\arg p$ 的平移与 $\beta(t_a)-\beta(t)$ 一致到 $10^{-16}$。
\end{proof}

\begin{proposition}[合成律与非交换余项]\status{已证+数值}\label{prop:comp}
两步合成满足
\begin{equation}
(p_2,q_2)\circ(p_1,q_1)=\big(p_2p_1+q_2\bar q_1,\;p_2q_1+q_2\bar p_1\big).
\end{equation}
对块内各步 $(p_k,q_k)$，记累积系数 $(P_k,Q_k)$ 与累积 Beltrami 系数 $M_k=Q_k/P_k$，则
\begin{equation}
\Big|\arg P_N-\sum_{k=1}^N\arg p_k\Big|\;\le\;\sum_{k=1}^{N}\arcsin\!\big(|\mu_k|\,|M_{k-1}|\big),\qquad |\mu_kM_{k-1}|<1 .
\end{equation}
\end{proposition}
\begin{proof}
由 $P_k=p_kP_{k-1}\big(1+\mu_k\bar M_{k-1}\bar P_{k-1}/P_{k-1}\big)$，以及 $|z|<1$ 时 $|\arg(1+z)|\le\arcsin|z|$，逐步累加即得。在 246 次随机平面激波穿越试验中，该上界没有一次被违反。
\end{proof}

\begin{proposition}[光滑流的一致性]\status{已证}\label{prop:smooth}
若单步映射 $F=I+\Delta t\,\nabla u+O(\Delta t^2)$，$\nabla u=S+W$，则
\begin{equation}
p=1+\tfrac12(\theta+i\omega)\Delta t+O(\Delta t^2),\qquad q=\tfrac12\big[(S_{11}-S_{22})+2iS_{12}\big]\Delta t+O(\Delta t^2),
\end{equation}
所以 $\arg p_k=\tfrac12\omega\Delta t+O(\Delta t^2)$、$|\mu_k|=O(|S_{\mathrm{dev}}|\Delta t)$。由命题~\ref{prop:comp}，在固定物质三角形上 $\sum_k\arg p_k=\tfrac12\int\omega\,dt+O(\Delta t)$（$\omega$ 取三角形平均，空间误差 $O(h)$）。
\end{proposition}

\begin{proposition}[单轴压缩的签名]\status{已证+数值}\label{prop:uni}
平面激波的完整穿越是单轴压缩 $F=I-(1-1/\eta)\,nn^\top$，满足
\begin{equation}
J=1/\eta,\qquad \arg p=0,\qquad |\mu|=\frac{\eta-1}{\eta+1}=\tanh\!\big(\tfrac12|\ln J|\big),\qquad r\equiv\ln J-(\mathrm{tr}\,F-2)=1-\tfrac1\eta-\ln\eta .
\end{equation}
\end{proposition}

各向同性压缩满足 $|\mu|=0$，纯剪切满足 $\ln J=0$，单轴压缩落在曲线 $|\mu|=\tanh(\frac12|\ln J|)$ 上（图~\ref{fig:checks}c）。这三种情况在 $(-\ln J,|\mu|)$ 平面上彼此分开。
%=====================================================================
\section{激波签名与锚定规则}\label{sec:anchor}

\begin{definition}[热步与合并穿越块]
对每个物质三角形逐步计算 $(p_k,q_k)$。若 $|\mu_k|>\mu^*$ 称第 $k$ 步为热步。阈值取背景的稳健倍数 $\mu^*=\max\big(\kappa\cdot\mathrm{median}|\mu|,\ \mu_{\min}\big)$，中位数在全部三角形与全部步上计算，$\kappa\approx5$–10。连续的热步（允许中间夹 $h_s$ 个冷步，$h_s=2$）用合成律并成一个块；冷步各自成块。
\end{definition}

由命题~\ref{prop:smooth}，光滑流中 $|\mu_k|=O(|S_{\mathrm{dev}}|\Delta t)$，所以阈值随 $\Delta t$ 自动缩放。激波穿越时 $|\mu_k|=O\big((\eta-1)/(\eta+1)\big)$，与 $\Delta t$ 无关。这与 C 通道依靠“面积台阶与 $\Delta t$ 无关”的逻辑相同，但 $|\mu|$ 还能看见面积不变的非仿射步（例如切向速度跳变跨过三角形）。

\begin{definition}[块的激波签名]
对每个块记签名 $\big(-\ln J_B,\ |\mu_B|,\ \arg p_B\big)$。若 $\big||\mu_B|-\tanh(\tfrac12|\ln J_B|)\big|\le\epsilon_\mu$ 且 $|\arg p_B|\le\epsilon_\psi$，判为平面激波块，压缩比 $\eta=e^{-\ln J_B}$。偏离这条曲线的量度量弯曲激波、剪切或斜压转动的贡献。
\end{definition}

\begin{proposition}[合并规则消除平面激波的伪转动]\status{数值}\label{prop:merge}
在平面激波扫过随机物质三角形的试验中（$\eta=1.5$–6，$\Delta t$ 取三档），逐步求和 $\sum\arg p_k$ 的伪转角中位数为 0.006–0.07 rad、95 分位为 0.02–0.20 rad，与爆炸波中观察到的 0.1–0.18 rad 同量级；按合并规则（$\mu^*=0.01$，$h_s=2$）计算，95 分位降到 $1\times10^{-3}$–$5\times10^{-3}$ rad。完整穿越块的 $\arg p_B$ 精确为零（$10^{-15}$），签名精确落在单轴曲线上。
\end{proposition}

\begin{figure}[h]
\centering
\includegraphics[width=\linewidth]{fig_checks.pdf}
\caption{命题的数值检验。(a) 平面激波扫过物质三角形时的伪转角：逐步求和与合并穿越块。(b) 0.05$h$ 位置噪声下转角误差随步数的累积：固定物质三角形近似望远相消，每步独立取样按 $\sqrt n$ 增长。(c) 激波签名平面：完整穿越块落在单轴曲线 $|\mu|=\tanh(\frac12|\ln J|)$ 上，光滑流单步聚在原点附近。}
\label{fig:checks}
\end{figure}

\begin{finding}
\keyp{统一锚定规则。} 账本 $\ln J$、转角 $\varphi$、环量都在同一组块映射上计算：块内用合成律，块边界处若 $\Lambda_t<0.2$ 或最大内角 $>0.86\pi$ 才重锚，且不在热步中间重锚。这取代了“由 C 通道推迟重锚”，C 通道由前提降为报告量：节点 $i$ 的激波时刻就是其星内热块的并集。
\end{finding}

\begin{remark}
Dirac 残差在块上写作 $r=\ln(|p|^2-|q|^2)-2(\re p-1)$。光滑流中为 $O(\Delta t^2)$；平面激波满足命题~\ref{prop:uni} 的 $r(\eta)$。它与签名 $(\ln J,|\mu|)$ 冗余，保留作光滑区的一致性检查。
\end{remark}

%=====================================================================
\section{离散相对转角与离散 LAVD}\label{sec:rot}

\begin{definition}[离散相对转角与离散 LAVD]
对三角形 $t$ 的第 $k$ 个块取 $\psi_t^k=\arg p_t^{B_k}$；节点值 $\psi_i^k$ 取星内三角形按面积加权的平均；区域均值 $\bar\psi^k=\sum_{t\subset U}A_t\psi_t^k\big/\sum_{t\subset U}A_t$。定义
\begin{equation}
\varphi_i=\sum_k\big(\psi_i^k-\bar\psi^k\big),\qquad
\mathrm{LAVD}_i=2\sum_k\big|\psi_i^k-\bar\psi^k\big|-b_i ,
\end{equation}
其中 $b_i$ 是噪底（命题~\ref{prop:noise}）。区域均值逐步计算，只用当步为冷步的三角形；热块覆盖多步时，扣除量取该块所含各步区域均值之和。块的平移等于各步平移之和，所以严格客观性不受影响。
\end{definition}

\begin{proposition}[严格客观]\status{已证}\label{prop:obj}
在 $x=Q(t)y+b(t)$ 下，$\varphi_i$ 与 $\mathrm{LAVD}_i$ 严格不变。
\end{proposition}
\begin{proof}
由命题~\ref{prop:frame}，同一块上所有三角形的 $\psi_t^k$ 都平移 $-\Delta\beta_k$。面积加权的节点值与区域均值同样平移 $-\Delta\beta_k$，差值不变。面积本身在刚体变换下不变。
\end{proof}

此前实验用的是 $c_t/A_t=2\im p$，它在旋转参考系下的残差为 $-\sin\alpha\,(\theta_t-\bar\theta)\Delta t$。在可压缩流中这一项不均匀，所以那个版本只客观到 $O(\Delta t)$。改用 $\arg p$ 后客观性是严格的。

\begin{proposition}[一致性]\status{已证}\label{prop:cons}
在光滑区域内，$\varphi_i\to\tfrac12\int(\omega-\bar\omega)\,dt$（Haller 2016 的二维相对动态转角，$\bar\omega$ 为 $U$ 上的面积加权均值），$\mathrm{LAVD}_i+b_i\to\int|\omega-\bar\omega|\,dt$。误差为 $O(\Delta t)+O(h)$。
\end{proposition}
\begin{proof}
由命题~\ref{prop:smooth}，$\psi_t^k=\tfrac12\langle\omega\rangle_t\Delta t_k+O(\Delta t^2)$；面积加权平均把三角形平均转为节点值与区域均值，空间误差 $O(h)$。冷步上块长为一步，热步只占有限步数，不影响阶数。
\end{proof}

\begin{proposition}[噪声的累积方式]\status{数值}\label{prop:noise}
位置噪声 $\sigma$ 下，固定物质三角形上 $\sum_k\arg p_k$ 的误差近似望远相消：0.05$h$ 噪声下，$n=10,20,40,80,160$ 步的标准差为 0.070、0.071、0.075、0.078、0.082 rad；而每步独立取样时为 0.22、0.31、0.42、0.60、0.86 rad，按 $\sqrt n$ 增长（图~\ref{fig:checks}b）。所以带符号的 $\varphi$ 应在物质三角形上锚定到形状门失效为止，逐帧重剖是现有实验中噪声的主要来源。$|\cdot|$ 不能望远相消，LAVD 的噪底取
\begin{equation}
b_i=2\,n_B\,s_i\sqrt{2/\pi},
\end{equation}
$s_i$ 由块转角序列二阶差分的 MAD 估计，$n_B$ 为块数。
\end{proposition}
\begin{proof}[依据]
$\arg p_k$ 在小形变下近似等于三角形取向角 $\vartheta_{k+1}-\vartheta_k$ 加上形变与噪声的交叉项，求和后只剩端点噪声 $O(\sigma/h)$ 与交叉项的弱增长。每步独立取样破坏了这种相消。
\end{proof}

\begin{proposition}[激波]\status{已证+数值}
在合并规则下，平面无旋激波对 $\varphi$ 的贡献为完整穿越块的 $\arg p_B$，精确为零；剩余误差来自未被并入热块的边缘步，数值上 $\le5\times10^{-3}$ rad。弯曲激波上，块的极转角给出物质三角形在穿越期间的净刚体转动，与 Crocco 涡量沉积相对应；其精度需在激波–气泡 DNS 上检验。
\end{proposition}

\begin{definition}[涡边界]
在锚定构形（初始规则播种的三角剖分）上，取包围 $\varphi$ 或 LAVD 局部极大的最外层闭合等值线，要求凸性亏损 $\le d^*$（Haller 等取 $10^{-3}$）且弧长不小于数据可分辨的最小涡周长。不凸的涡（激波–气泡涡环、$\lambda$ 脚涡）需对 $d^*$ 做敏感性扫描；第~\ref{sec:persist} 节给出不依赖凸性的替代。
\end{definition}

%=====================================================================
\section{物质账本与环量收支}

\begin{proposition}[块链账本]\status{已证}
节点账本 $\ln J_i=\sum_k\ln J^{B_k}_i$，其中星面积比在块映射上计算。平面激波前后两态均匀时，物质三角形的边保持直线，完整穿越块的面积比精确等于 $1/\eta$，跨激波的面积跳变精确为 $-\ln\eta$。
\end{proposition}

这条命题与此前“由 C 通道推迟重锚使望远相消”的结论相同，只是触发条件换成了三角形自身的 $|\mu|$，因此不需要全局检测。环量部分保持已有结论：区域环量是边界线积分，与内部剖分无关；物质三角形的环量跳变 $\Delta\Gamma_t/A_t(t_a)$ 取块端点三步位移环量的中位数，替代逐点力矩图。

%=====================================================================
\section{物质胞腔复形与空心胞腔}\label{sec:cells}

\begin{definition}[物质胞腔复形]
块边界处若某三角形未通过形状门，不在局部重新剖分，而是把它与共享其最长边的邻居合并为物质多边形，重复直到等周商 $4\pi A_P/P_P^2\ge\Lambda^*$。2-胞腔为多边形，边界矩阵 $B_2$ 由多边形的有向边界给出。
\end{definition}

\begin{proposition}[可加性]\status{已证}
$c_P=\sum_{t\subset P}c_t$（内部边相互抵消），$A_P=\sum_{t\subset P}A_t$。面积比由多边形顶点直接计算，精确；$c_P/A_P$ 等于各三角形 $2\im p_t$ 的面积加权平均。环量跳变与 Kelvin 收支在胞腔上照常成立。
\end{proposition}

\begin{proposition}[噪声随合并降低]\status{量级}
位置噪声 $\sigma$ 下，$\delta(c_P/A_P)\sim\sigma P_P/A_P$。由 $m$ 个尺寸为 $h$ 的三角形合并成紧凑胞腔时，$P_P\propto\sqrt m\,h$、$A_P\propto mh^2$，误差按 $m^{-1/2}$ 下降。胞腔合并因此同时缓解“重锚放大噪声、不重锚积累几何误差”的取舍。
\end{proposition}

\begin{definition}[非仿射度与撕裂]
对顶点数不少于 4 的胞腔 $P$，定义
\begin{equation}
\varepsilon_P=\frac{1}{\mathrm{diam}\,P}\min_{F,c}\Big[\frac1{n_v}\sum_{v\in\partial P}|x_v-FX_v-c|^2\Big]^{1/2}.
\end{equation}
若 $\varepsilon_P>\varepsilon^*=\max\big(c_\sigma\sigma/\mathrm{diam}\,P,\ \varepsilon_0\big)$，判该胞腔被撕裂，保持空心（不填充）。物理含义是：这块物质已不再作为一个仿射体运动，流动的非线性在胞腔尺度上不可忽略（拉丝、卷起、滑移线）。数据驱动的替代判据是 Sardellitti–Barbarossa（2024）的“边信号沿胞腔的总变差加稀疏”准则。
\end{definition}

\begin{proposition}[空心胞腔与调和空间]\status{已证}
若全部胞腔填充后的复形同胚于圆盘，则留空 $k$ 个胞腔后 $\beta_1=k$；对任意对称正定的度量矩阵，调和 1-上链空间的维数都等于 $k$。
\end{proposition}
\begin{proof}
Euler 示性数 $\chi=V-E+(F-k)=1-k$；平面复形 $\beta_2=0$，连通时 $\beta_0=1$，故 $\beta_1=k$。Hodge 定理对任意内积成立，$\dim\ker L_1=\beta_1$。
\end{proof}

调和流绕空心胞腔环流，所以撕裂的位置与持续时间可以直接读出，可作为屏障与拉丝的拓扑指示，不需要人为阈值之外的参数。

%=====================================================================
\section{Hodge–Dirac 层的一般内积形式}

\begin{proposition}[一般内积下的 Hodge 分解]\status{已证}
在任意满足 $d_1d_0=0$ 的胞腔复形上，对任意对称正定的 $M_0,M_1,M_2$，有 $C^1=\mathrm{im}\,d_0\oplus\mathrm{im}\,\delta_2\oplus\ker L_1$，三部分在 $\langle\cdot,\cdot\rangle_{M_1}$ 下两两正交，$\delta_k=M_{k-1}^{-1}d_{k-1}^\top M_k$。
\end{proposition}

这就是 Lim（2020）在图上的框架推广到胞腔复形的形式，可以同时解决 T-08 冲突：T3–T5 的命题全部改写为“任意对称正定内积”，对角 cotan 只作为特例。推荐的度量为：三角形取 Whitney/Galerkin 质量（形状正则即收敛）；多边形取模仿有限差分或虚单元的内积（对线性场一致、对网格稳定）。组合 Laplacian（AFRC 所用）不含 Hodge 星，不给出物理分解（Ribando-Gros 等，2024），所以凡是涉及份额、谱、扩散的命题一律放在带 $M_k$ 的复形上。

信号取块上的位移 1-上链 $f_e^{B}=\tfrac12(\Delta x_a+\Delta x_b)\cdot(x_b-x_a)$。$D(0,f,0)=(\delta_1f,0,d_1f)$ 的节点分量是弱散度，三角形或胞腔分量是位移环量。只报告在块或窗口上积分后的份额，不报告逐点散度或涡量。扇区分离与“线性声学即无质量 Dirac 方程”两条结论不变。

%=====================================================================
\section{阈值消除：持续拓扑}\label{sec:persist}

\begin{definition}[过滤]
在锚定胞腔复形上做 $|\varphi|$ 的上水平集过滤：$K_w$ 由全部顶点满足 $|\varphi|\ge w$ 的胞腔组成。长寿命的 $\beta_0$ 类为涡核。屏障一侧用 $\varepsilon_P$ 的过滤：按 $\varepsilon_P$ 从大到小把胞腔挖空，长寿命的 $\beta_1$ 类为持续的撕裂。
\end{definition}

\begin{proposition}[稳定性]\status{引用}
持续图对过滤函数是 1-Lipschitz 的：$d_B(\mathrm{Dgm}\,f,\mathrm{Dgm}\,g)\le\|f-g\|_\infty$（Cohen-Steiner、Edelsbrunner、Harer 2007）。
\end{proposition}

结合命题~\ref{prop:noise}：固定锚定下 $\varphi$ 的噪声水平 $\epsilon_\varphi$ 由端点噪声决定、不随步数增长，所以“持续性大于 $2\epsilon_\varphi$ 的核”是由噪声推出的判据，取代 top-10\%、$\pi/4$ 和分水岭合并距离。多个窗口串成藤图后，涡的出生、合并与死亡就是 $\beta_0$ 类的出生与死亡；涡与间隙鞍点同时消失满足离散 Morse 恒等式 $\#\max-\#\mathrm{saddle}+\#\min=\chi$。

%=====================================================================
\section{与输运、掺混层的接口}

\begin{itemize}
\item \keyp{扩散网格与运动学网格分开。} 网络扩散的离散极大值原理当且仅当每条内部边的对角和 $\le\pi$（$\cot\alpha+\cot\beta\ge0$）时成立，所以扩散用当前 Delaunay 剖分，在 Delaunay 条件失效前重剖。三角形层的运动学量用锚定的物质三角形与胞腔。两者共用示踪点与质量。
\item \keyp{切向拉伸的复形式。} 单位切向 $T=e^{i\vartheta_T}$ 时 $|FT|=|p+q\,e^{-2i\vartheta_T}|$。沿块用合成律串联 $(p,q)$，得到“三点链式拉伸”：对分片仿射运动精确，不需要 12 近邻最小二乘，噪声由命题~\ref{prop:noise} 的锚定规则控制。掺混分解 $\ln|\nabla c|/|\nabla c_0|=\ln\lambda-\ln J$ 的两项都由同一组 $(p,q)$ 给出。
\item \keyp{转移算子、质量通量与片层钟} 不受本稿影响。
\end{itemize}

%=====================================================================
\section{状态与最小检验}

\begin{center}\small
\begin{xltabular}{\linewidth}{@{}P{4.3cm}P{1.9cm}LP{2.6cm}@{}}
\toprule
命题 & 状态 & 最小检验（六算例） & 通过标准\\\midrule
复表示恒等式、参考系变换、合成律 & 已证+数值 & 已完成 & 机器精度\\
合并规则消除平面激波伪转动 & 数值 & Sod 与柱面爆炸波：$\varphi$ 的残余 & Sod $<10^{-3}$ rad；爆炸波从 0.1–0.18 降到 $<0.01$ rad\\
激波签名与压缩比 & 已证+数值 & Sod、爆炸波、激波–气泡：块签名与单轴曲线的偏离 & 平面激波处偏离 $<0.01$；气泡界面处偏离与 DNS 斜压项同号\\
$\varphi$ 严格客观、收敛到 Haller 转角 & 已证 & 双环流与涡对：旋转参考系与网格加密 & 客观性误差在机器精度；与解析 $\varphi$ 的 Spearman $\ge0.99$\\
固定锚定下噪声望远相消 & 数值 & 激波–气泡加 0.1$h$ 噪声：锚定 $\varphi$ 对 DNS & Spearman 由 0.59 提高（目标 $\ge0.8$）\\
LAVD 噪底 & 待检验 & 涡对加噪声 & 与解析 LAVD 相关 $\ge0.9$\\
胞腔合并降噪 & 量级 & 激波–气泡加 0.1$h$ 噪声：环量跳变 & 相关由 0.41 提高到 $\ge0.8$\\
空心胞腔与 $\beta_1$ & 已证 & 激波–气泡：撕裂位置与 DNS 卷起 & 位置一致，且对 $\varepsilon^*$ 稳健\\
持续性核数 & 引用+待检验 & 涡对：$\beta_0$ 持续图 & 2 个核的持续性 $>2\epsilon_\varphi$，其余 $<2\epsilon_\varphi$\\
三点链式拉伸 & 已证 & 六算例界面拉伸对参照 & 不劣于 12 近邻链式张量\\
\bottomrule
\end{xltabular}\end{center}

尚未补上的理论空白有三处：弯曲激波上块极转角对斜压沉积的精度，需要激波–气泡 DNS 来定；多边形胞腔上内积的具体构造，可以直接采用模仿有限差分的现成形式；三维推广没有复表示可用，需改用四面体映射的极分解，Beltrami 系数换成奇异值比。

\section*{参考文献}
\addcontentsline{toc}{section}{参考文献}
{\small
\begin{enumerate}[label={[\arabic*]},leftmargin=2em,itemsep=0.1em]
\item G. Haller, A. Hadjighasem, M. Farazmand, F. Huhn, Defining coherent vortices objectively from the vorticity, J.\ Fluid Mech.\ 795, 136–173 (2016).
\item G. Haller, Dynamic rotation and stretch tensors from a dynamic polar decomposition, J.\ Mech.\ Phys.\ Solids 86, 70–93 (2016).
\item L.-H. Lim, Hodge Laplacians on graphs, SIAM Rev.\ 62, 685–715 (2020).
\item E. Ribando-Gros, R. Wang, J. Chen, Y. Tong, G.-W. Wei, Combinatorial and Hodge Laplacians: similarities and differences, SIAM Rev.\ (2024).
\item S. Sardellitti, S. Barbarossa, Topological signal processing over generalized cell complexes, IEEE Trans.\ Signal Process.\ (2024).
\item T. M. Roddenberry, M. T. Schaub, M. Hajij, Signal processing on cell complexes, ICASSP (2022).
\item D. Cohen-Steiner, H. Edelsbrunner, J. Harer, Stability of persistence diagrams, Discrete Comput.\ Geom.\ 37, 103–120 (2007).
\item C. Zhu, S. H. Christiansen, K. Hu, A. N. Hirani, Convergence and stability of discrete exterior calculus for the Hodge Laplace problem in two dimensions, arXiv:2505.08966 (2026).
\item H. S. Huntley et al., Conditions for reliable divergence estimates from drifter triplets, J.\ Atmos.\ Oceanic Technol.\ 39, 1499 (2022).
\item L. V. Ahlfors, Lectures on Quasiconformal Mappings（Beltrami 系数与合成律的标准来源）.
\end{enumerate}
}
\end{document}
